% Options for packages loaded elsewhere
\PassOptionsToPackage{unicode}{hyperref}
\PassOptionsToPackage{hyphens}{url}
\documentclass[
]{article}
\usepackage{xcolor}
\usepackage[margin=1in]{geometry}
\usepackage{amsmath,amssymb}
\setcounter{secnumdepth}{-\maxdimen} % remove section numbering
\usepackage{iftex}
\ifPDFTeX
  \usepackage[T1]{fontenc}
  \usepackage[utf8]{inputenc}
  \usepackage{textcomp} % provide euro and other symbols
\else % if luatex or xetex
  \usepackage{unicode-math} % this also loads fontspec
  \defaultfontfeatures{Scale=MatchLowercase}
  \defaultfontfeatures[\rmfamily]{Ligatures=TeX,Scale=1}
\fi
\usepackage{lmodern}
\ifPDFTeX\else
  % xetex/luatex font selection
\fi
% Use upquote if available, for straight quotes in verbatim environments
\IfFileExists{upquote.sty}{\usepackage{upquote}}{}
\IfFileExists{microtype.sty}{% use microtype if available
  \usepackage[]{microtype}
  \UseMicrotypeSet[protrusion]{basicmath} % disable protrusion for tt fonts
}{}
\makeatletter
\@ifundefined{KOMAClassName}{% if non-KOMA class
  \IfFileExists{parskip.sty}{%
    \usepackage{parskip}
  }{% else
    \setlength{\parindent}{0pt}
    \setlength{\parskip}{6pt plus 2pt minus 1pt}}
}{% if KOMA class
  \KOMAoptions{parskip=half}}
\makeatother
\usepackage{color}
\usepackage{fancyvrb}
\newcommand{\VerbBar}{|}
\newcommand{\VERB}{\Verb[commandchars=\\\{\}]}
\DefineVerbatimEnvironment{Highlighting}{Verbatim}{commandchars=\\\{\}}
% Add ',fontsize=\small' for more characters per line
\usepackage{framed}
\definecolor{shadecolor}{RGB}{248,248,248}
\newenvironment{Shaded}{\begin{snugshade}}{\end{snugshade}}
\newcommand{\AlertTok}[1]{\textcolor[rgb]{0.94,0.16,0.16}{#1}}
\newcommand{\AnnotationTok}[1]{\textcolor[rgb]{0.56,0.35,0.01}{\textbf{\textit{#1}}}}
\newcommand{\AttributeTok}[1]{\textcolor[rgb]{0.13,0.29,0.53}{#1}}
\newcommand{\BaseNTok}[1]{\textcolor[rgb]{0.00,0.00,0.81}{#1}}
\newcommand{\BuiltInTok}[1]{#1}
\newcommand{\CharTok}[1]{\textcolor[rgb]{0.31,0.60,0.02}{#1}}
\newcommand{\CommentTok}[1]{\textcolor[rgb]{0.56,0.35,0.01}{\textit{#1}}}
\newcommand{\CommentVarTok}[1]{\textcolor[rgb]{0.56,0.35,0.01}{\textbf{\textit{#1}}}}
\newcommand{\ConstantTok}[1]{\textcolor[rgb]{0.56,0.35,0.01}{#1}}
\newcommand{\ControlFlowTok}[1]{\textcolor[rgb]{0.13,0.29,0.53}{\textbf{#1}}}
\newcommand{\DataTypeTok}[1]{\textcolor[rgb]{0.13,0.29,0.53}{#1}}
\newcommand{\DecValTok}[1]{\textcolor[rgb]{0.00,0.00,0.81}{#1}}
\newcommand{\DocumentationTok}[1]{\textcolor[rgb]{0.56,0.35,0.01}{\textbf{\textit{#1}}}}
\newcommand{\ErrorTok}[1]{\textcolor[rgb]{0.64,0.00,0.00}{\textbf{#1}}}
\newcommand{\ExtensionTok}[1]{#1}
\newcommand{\FloatTok}[1]{\textcolor[rgb]{0.00,0.00,0.81}{#1}}
\newcommand{\FunctionTok}[1]{\textcolor[rgb]{0.13,0.29,0.53}{\textbf{#1}}}
\newcommand{\ImportTok}[1]{#1}
\newcommand{\InformationTok}[1]{\textcolor[rgb]{0.56,0.35,0.01}{\textbf{\textit{#1}}}}
\newcommand{\KeywordTok}[1]{\textcolor[rgb]{0.13,0.29,0.53}{\textbf{#1}}}
\newcommand{\NormalTok}[1]{#1}
\newcommand{\OperatorTok}[1]{\textcolor[rgb]{0.81,0.36,0.00}{\textbf{#1}}}
\newcommand{\OtherTok}[1]{\textcolor[rgb]{0.56,0.35,0.01}{#1}}
\newcommand{\PreprocessorTok}[1]{\textcolor[rgb]{0.56,0.35,0.01}{\textit{#1}}}
\newcommand{\RegionMarkerTok}[1]{#1}
\newcommand{\SpecialCharTok}[1]{\textcolor[rgb]{0.81,0.36,0.00}{\textbf{#1}}}
\newcommand{\SpecialStringTok}[1]{\textcolor[rgb]{0.31,0.60,0.02}{#1}}
\newcommand{\StringTok}[1]{\textcolor[rgb]{0.31,0.60,0.02}{#1}}
\newcommand{\VariableTok}[1]{\textcolor[rgb]{0.00,0.00,0.00}{#1}}
\newcommand{\VerbatimStringTok}[1]{\textcolor[rgb]{0.31,0.60,0.02}{#1}}
\newcommand{\WarningTok}[1]{\textcolor[rgb]{0.56,0.35,0.01}{\textbf{\textit{#1}}}}
\usepackage{graphicx}
\makeatletter
\newsavebox\pandoc@box
\newcommand*\pandocbounded[1]{% scales image to fit in text height/width
  \sbox\pandoc@box{#1}%
  \Gscale@div\@tempa{\textheight}{\dimexpr\ht\pandoc@box+\dp\pandoc@box\relax}%
  \Gscale@div\@tempb{\linewidth}{\wd\pandoc@box}%
  \ifdim\@tempb\p@<\@tempa\p@\let\@tempa\@tempb\fi% select the smaller of both
  \ifdim\@tempa\p@<\p@\scalebox{\@tempa}{\usebox\pandoc@box}%
  \else\usebox{\pandoc@box}%
  \fi%
}
% Set default figure placement to htbp
\def\fps@figure{htbp}
\makeatother
\setlength{\emergencystretch}{3em} % prevent overfull lines
\providecommand{\tightlist}{%
  \setlength{\itemsep}{0pt}\setlength{\parskip}{0pt}}
\usepackage{bookmark}
\IfFileExists{xurl.sty}{\usepackage{xurl}}{} % add URL line breaks if available
\urlstyle{same}
\hypersetup{
  pdftitle={Gibbs Sampling},
  hidelinks,
  pdfcreator={LaTeX via pandoc}}

\title{Gibbs Sampling}
\author{}
\date{\vspace{-2.5em}2026-04-17}

\begin{document}
\maketitle

\subsection{Introduction}\label{introduction}

Gibbs sampling is one of the oldest practical MCMC algorithms and
remains an instructive special case of more general samplers. It
exploits a structural property of multidimensional posteriors: even when
the joint distribution is intractable, the conditional distribution of
one parameter given all others can often be derived in closed form.
Cycling through the parameters and updating each from its full
conditional generates a Markov chain whose stationary distribution is
the joint posterior, with no acceptance step and no tuning of proposals.

\subsection{Prerequisites}\label{prerequisites}

A working understanding of MCMC, conditional probability, and the
conjugate-prior structure that makes full conditionals tractable for
canonical models (Normal--Normal, Beta--Binomial, Gamma--Poisson,
Dirichlet--Multinomial).

\subsection{Theory}\label{theory}

For a joint posterior \(p(\theta_1, \dots, \theta_K \mid y)\), the
algorithm initialises at some \(\theta^{(0)}\) and, at iteration \(t\),
sweeps through the parameters in order:

\[\theta_k^{(t)} \sim p\bigl(\theta_k \mid \theta_1^{(t)}, \dots, \theta_{k-1}^{(t)}, \theta_{k+1}^{(t-1)}, \dots, \theta_K^{(t-1)}, y\bigr)\]

for \(k = 1, \dots, K\). Under standard regularity conditions, the chain
is irreducible and aperiodic and converges to the joint posterior. Each
update is a draw from a one-dimensional or low-dimensional conditional,
often via standard distribution libraries; no Metropolis acceptance is
needed because the proposal \emph{is} the conditional. When some
conditionals are not available analytically, the affected blocks are
updated with a Metropolis step embedded inside the Gibbs sweep
(``Metropolis-within-Gibbs'').

\subsection{Assumptions}\label{assumptions}

Every full conditional must be either a recognisable distribution or a
distribution one can sample from with a Metropolis step. Strongly
correlated parameters update slowly under a coordinate-wise sweep; this
is the central pathology of Gibbs and the reason HMC is now preferred
for general continuous models.

\subsection{R Implementation}\label{r-implementation}

\begin{Shaded}
\begin{Highlighting}[]
\FunctionTok{set.seed}\NormalTok{(}\DecValTok{2026}\NormalTok{)}

\CommentTok{\# Conjugate normal{-}normal model: mu and sigma2 both unknown}
\CommentTok{\# Priors: mu \textasciitilde{} N(mu0, tau0\^{}2); 1/sigma2 \textasciitilde{} Gamma(alpha, beta)}
\NormalTok{y }\OtherTok{\textless{}{-}} \FunctionTok{rnorm}\NormalTok{(}\DecValTok{50}\NormalTok{, }\DecValTok{3}\NormalTok{, }\DecValTok{2}\NormalTok{)}
\NormalTok{n }\OtherTok{\textless{}{-}} \FunctionTok{length}\NormalTok{(y)}
\NormalTok{mu0 }\OtherTok{\textless{}{-}} \DecValTok{0}\NormalTok{; tau0 }\OtherTok{\textless{}{-}} \DecValTok{10}
\NormalTok{alpha }\OtherTok{\textless{}{-}} \DecValTok{2}\NormalTok{; beta }\OtherTok{\textless{}{-}} \DecValTok{1}

\NormalTok{n\_iter }\OtherTok{\textless{}{-}} \DecValTok{5000}
\NormalTok{mu\_chain    }\OtherTok{\textless{}{-}} \FunctionTok{numeric}\NormalTok{(n\_iter)}
\NormalTok{sig2\_chain  }\OtherTok{\textless{}{-}} \FunctionTok{numeric}\NormalTok{(n\_iter)}
\NormalTok{sig2\_chain[}\DecValTok{1}\NormalTok{] }\OtherTok{\textless{}{-}} \FunctionTok{var}\NormalTok{(y)}

\ControlFlowTok{for}\NormalTok{ (t }\ControlFlowTok{in} \DecValTok{2}\SpecialCharTok{:}\NormalTok{n\_iter) \{}
  \CommentTok{\# Sample mu | sigma\^{}2, y}
\NormalTok{  prec }\OtherTok{\textless{}{-}} \DecValTok{1}\SpecialCharTok{/}\NormalTok{tau0}\SpecialCharTok{\^{}}\DecValTok{2} \SpecialCharTok{+}\NormalTok{ n}\SpecialCharTok{/}\NormalTok{sig2\_chain[t}\DecValTok{{-}1}\NormalTok{]}
\NormalTok{  mean\_mu }\OtherTok{\textless{}{-}}\NormalTok{ (mu0}\SpecialCharTok{/}\NormalTok{tau0}\SpecialCharTok{\^{}}\DecValTok{2} \SpecialCharTok{+} \FunctionTok{sum}\NormalTok{(y)}\SpecialCharTok{/}\NormalTok{sig2\_chain[t}\DecValTok{{-}1}\NormalTok{]) }\SpecialCharTok{/}\NormalTok{ prec}
\NormalTok{  mu\_chain[t] }\OtherTok{\textless{}{-}} \FunctionTok{rnorm}\NormalTok{(}\DecValTok{1}\NormalTok{, mean\_mu, }\FunctionTok{sqrt}\NormalTok{(}\DecValTok{1}\SpecialCharTok{/}\NormalTok{prec))}

  \CommentTok{\# Sample sigma\^{}2 | mu, y  (inverse gamma)}
\NormalTok{  alpha\_n }\OtherTok{\textless{}{-}}\NormalTok{ alpha }\SpecialCharTok{+}\NormalTok{ n}\SpecialCharTok{/}\DecValTok{2}
\NormalTok{  beta\_n  }\OtherTok{\textless{}{-}}\NormalTok{ beta }\SpecialCharTok{+} \FunctionTok{sum}\NormalTok{((y }\SpecialCharTok{{-}}\NormalTok{ mu\_chain[t])}\SpecialCharTok{\^{}}\DecValTok{2}\NormalTok{)}\SpecialCharTok{/}\DecValTok{2}
\NormalTok{  sig2\_chain[t] }\OtherTok{\textless{}{-}} \DecValTok{1}\SpecialCharTok{/}\FunctionTok{rgamma}\NormalTok{(}\DecValTok{1}\NormalTok{, alpha\_n, beta\_n)}
\NormalTok{\}}

\CommentTok{\# Summaries}
\FunctionTok{c}\NormalTok{(}\AttributeTok{mean\_mu =} \FunctionTok{mean}\NormalTok{(mu\_chain[}\SpecialCharTok{{-}}\NormalTok{(}\DecValTok{1}\SpecialCharTok{:}\DecValTok{1000}\NormalTok{)]),}
  \AttributeTok{mean\_sigma =} \FunctionTok{mean}\NormalTok{(}\FunctionTok{sqrt}\NormalTok{(sig2\_chain[}\SpecialCharTok{{-}}\NormalTok{(}\DecValTok{1}\SpecialCharTok{:}\DecValTok{1000}\NormalTok{)])))}
\end{Highlighting}
\end{Shaded}

\subsection{Output \& Results}\label{output-results}

The chain produces 5,000 draws of \(\mu\) and \(\sigma^2\). Discarding
the first 1,000 as burn-in and summarising the rest gives posterior
means for \(\mu\) near 3.1 and for \(\sigma\) near 2.0 --- close to the
data-generating parameters and within the expected Monte Carlo error.

\subsection{Interpretation}\label{interpretation}

A reporting sentence: ``A Gibbs sampler for the conjugate
Normal--Inverse-Gamma model produced 4,000 post-burn-in draws with
posterior mean \(\mu = 3.10\) (95 \% credible interval 2.55 to 3.65) and
posterior mean \(\sigma = 2.01\) (95 \% CrI 1.65 to 2.49).'' For
pedagogical examples, also report the rejection rate (Gibbs has none)
and the autocorrelation function so readers can see the per-iteration
efficiency.

\subsection{Practical Tips}\label{practical-tips}

\begin{itemize}
\tightlist
\item
  Block-update correlated parameters jointly when the conditional is
  available in closed form; this reduces autocorrelation dramatically
  compared with coordinate-wise updates.
\item
  For non-conjugate components, embed a Metropolis step
  (Metropolis-within-Gibbs); JAGS does this automatically when it cannot
  resolve a node analytically.
\item
  Check convergence with \(\hat R\) across multiple chains and trace
  plots; a single chain cannot reveal mode-trapping.
\item
  For continuous, smooth posteriors of moderate-to-high dimension, HMC
  (Stan) typically mixes far better than Gibbs because it follows the
  gradient instead of cycling through coordinates.
\item
  Gibbs remains the algorithm of choice for discrete-state-space models,
  latent-allocation samplers, and any setting where conjugacy makes a
  single sweep almost free.
\end{itemize}

\subsection{R Packages Used}\label{r-packages-used}

Base R for hand-coded Gibbs, \texttt{rjags} and \texttt{nimble} for
declarative model specification with automatic Gibbs--Metropolis
dispatch, and \texttt{coda} for diagnostic plots and convergence
statistics on the resulting chains.

\subsection{For Reviewers}\label{for-reviewers}

\textbf{What to look for in a paper using this method.}

\begin{itemize}
\tightlist
\item
  \textbf{Common misapplications.}
\item
  \textbf{Diagnostics that should be reported but often aren't.}
\item
  \textbf{Red flags in tables and figures.}
\item
  \textbf{What to verify.}
\item
  \textbf{What an adequate Methods paragraph must contain.}
\end{itemize}

\subsection{See also --- labs in this
chapter}\label{see-also-labs-in-this-chapter}

\begin{itemize}
\tightlist
\item
  \href{labs/lab-bayesian-thinking-control-vs-decision-error.Rmd}{Bayesian
  thinking; α control vs decision error}
\item
  \href{labs/lab-brms-stan-loo-hierarchical-models.Rmd}{brms / Stan,
  LOO, hierarchical models}
\end{itemize}

\end{document}
